\section{Implementation in the BX3 Framework}

The preceding sections established the Bailout Protocol's formal semantics, correctness properties, and operational mechanics as a stand-alone accountability architecture. This section demonstrates how the protocol is instantiated within the BX3 Framework's three-layer model, showing that the Bailout Protocol is not an overlay or optional enhancement but rather the \emph{runtime expression} of BX3's upstream accountability guarantee—a guarantee that is structurally embedded in the framework's design rather than grafted onto it post hoc. We examine each layer's role in protocol execution, the Bounds Engine as a cross-layer monitoring mechanism, the Fact Layer's enforcement of the immutable audit trail, and the state machine implementation as a Fact Layer extension recording forensic ledger entries.

\subsection{Purpose Layer as Human Accountability Anchor}

The BX3 Framework's \textbf{Purpose Layer} corresponds to what earlier sections termed the \emph{Intention} layer—a normative encoding of the goals, boundary constraints, and accountability assignments that govern agent behavior before any execution begins. In BX3, the Purpose Layer is the site of the \emph{human accountability anchor}: the point at which a designated principal explicitly authorises an agent's operating envelope, defines the bounds of permissible action, and records the justification for those bounds. This encoding is not merely informational; it is operationally loaded by the Bounds Engine at session initialization and enforced throughout the session's lifetime.

Concretely, the Purpose Layer entry for a given agent session contains three canonical fields. The first is the \textbf{principal identifier} $p$: the authenticated human or institutional entity that bears ultimate accountability for the agent's actions. The second is the \textbf{boundary specification} $B$, a set of constraints on the agent's action space—domain restrictions, value constraints, escalation thresholds, and intervention windows—that together define what the agent is permitted to do without further authorization. The third is the \textbf{justification record} $J$, a natural-language or structured attestation explaining why these bounds were set as they were, which serves as the evidentiary basis for downstream audit. The triplet $(p, B, J)$ constitutes the Purpose Layer's canonical state for a session.

When the Bailout Protocol initializes, it ingests $(p, B, J)$ from the Purpose Layer and uses these values to parameterize its enforcement logic. The principal $p$ is inscribed in the session's attribution chain and serves as the recipient of all escalation events. The boundary specification $B$ is decomposed into a set of runtime monitors—timing constraints, action-class restrictions, resource ceilings—that the Bounds Engine evaluates continuously. The justification record $J$ is written verbatim into the first entry of the forensic ledger, establishing the evidentiary baseline against which all subsequent events are evaluated. This integration is structural: the Bailout Protocol cannot initialize without a Purpose Layer entry, ensuring that no agent session can enter an operational state without a named principal and an auditable justification.

The Purpose Layer's role as an accountability anchor has a critical implication for the Bailout Protocol's correctness claims. Properties S (Safety), L (Liveness), and A (Accountability) from Section~\ref{sec:guarantees} are predicated on the existence of a non-null principal $p$. If the Purpose Layer entry is absent or malformed, the Bailout Protocol detects this at initialization and refuses to enter an active state—transitioning instead to a \texttt{boot\underline{ }failure} state with a null attribution. This is a deliberate design choice: the protocol treats a missing accountability anchor as a protocol-level failure, not merely a configuration warning. The rationale is that an agent operating without a named principal is, by definition, unaccountable—and permitting unaccountable operation violates the framework's founding commitment to human authority.

\subsection{Bounds Engine: Cross-Layer Bail-Out Triggering}

The \textbf{Bounds Engine} is BX3's cross-layer monitoring subsystem, operating at the intersection of the Purpose Layer and the Execution/Fact Layer. Its function is to evaluate the agent's running state against the boundary specification $B$ sourced from the Purpose Layer, and to trigger bail-out events when the evaluated state violates or approaches violation of any constraint in $B$. The Bounds Engine is the primary mechanism by which the Bailout Protocol's Safety property (Property S) is enforced at runtime.

At initialization, the Bounds Engine receives a compiled form of $B$—the boundary specification—as a set of monotonic predicate functions over the agent's observable state variables. These predicates are evaluated at each checkpoint interval $\tau_{\text{ckpt}}$ against the current state snapshot $c_t$. Each predicate maps the state variables to a boolean verdict: \texttt{true} indicates the state is within bounds; \texttt{false} indicates a boundary violation. When any predicate returns \texttt{false}, the Bounds Engine does not wait for further evaluation cycles—it generates a \texttt{boundary\underline{ }violation} event and simultaneously: (i) writes the event to the forensic ledger with full state context, (ii) transitions the session to a \texttt{bailout\underline{ }pending} state, and (iii) notifies the principal $p$ via the registered escalation channel.

The Bounds Engine implements three categories of bounds monitoring, each corresponding to a different dimension of the boundary specification $B$:

\begin{itemize}
    \item \textbf{Temporal bounds} enforce that the elapsed time since the last authenticated human contact does not exceed the liveness threshold $T_{\text{liveness}}$. The temporal monitor evaluates the \texttt{last\underline{ }contact} timestamp at each checkpoint and generates a \texttt{timeout} event if the threshold is breached. This enforces Property L (Liveness) at the architectural level.
    \item \textbf{Spatial bounds} enforce that the agent's action space remains within the domain constraints specified in $B$. This includes action-class restrictions (e.g., \emph{may not send external communications without prior authorization}), resource ceilings (e.g., \emph{may not allocate more than $X$ compute units to a single task}), and object-space restrictions (e.g., \emph{may not access data outside designated namespaces}). Spatial bound violations are detected by the Fact Layer's state transition validator (Section~\ref{sec:fact-layer}) before they are committed to the forensic ledger.
    \item \textbf{Trajectory bounds} enforce that the agent's decision sequence does not diverge from the intended objective pathway by more than a configurable drift threshold. Trajectory monitoring is implemented as a rolling horizon comparison: the agent's current goal-state representation is compared against the authorized goal-state encoded in the Purpose Layer entry. When the cosine similarity or equivalent distance metric between the two representations falls below the drift threshold, a \texttt{drift} event is generated. Trajectory bounds are the mechanism by which the alignment property is enforced at runtime—they detect behavioral drift before it manifests as a concrete policy violation.
\end{itemize}

The Bounds Engine is designed to be failure-isolated from the agent it monitors. The engine runs as an independent process with its own execution context, message queue, and checkpoint schedule. An agent that becomes unresponsive or enters a runaway state does not affect the Bounds Engine's ability to evaluate predicates or generate bail-out events. This isolation is architecturally essential: a Bounds Engine co-located within the agent's process could itself be compromised by the agent's anomalous behavior, defeating the purpose of the external monitoring mechanism.

\subsection{Fact Layer: Audit Trail Enforcement}
\label{sec:fact-layer}

The \textbf{Fact Layer} is the BX3 layer responsible for recording the factual chain of events—the \emph{what happened} of agent execution. It corresponds to the \emph{Execution} layer in the framework's earlier formulation and is where the Bailout Protocol's forensic ledger lives. The Fact Layer enforces the audit trail through a strict append-only logging regime, a state transition validation protocol, and a cryptographic sealing mechanism that renders the ledger tamper-evident.

The Fact Layer's core data structure is the \textbf{forensic ledger entry} $e_t$, a tuple written at each checkpoint $c_t$ and at each significant state transition event:

\begin{equation}
e_t = \bigl(t, \sigma_t, a_t, \tau_t, \phi_t, p_t\bigr)
\end{equation}

where $t$ is the wall-clock timestamp, $\sigma_t$ is the session identifier, $a_t$ is the action or decision taken at this step, $\tau_t$ is the cryptographic seal (SHA-256 hash chained to the preceding entry), $\phi_t$ is the principal attribution chain, and $p_t$ is the Purpose Layer reference used to validate the action's boundary conformance. The chained hash $\tau_t$ is computed as:

\begin{equation}
\tau_t = H\bigl(\tau_{t-1} \parallel e_t\bigr)
\end{equation}

where $H$ is a cryptographic hash function and $\tau_{t-1}$ is the seal of the preceding entry. This creates an immutable hash chain: any modification of a historical entry breaks the chain at that point and all subsequent points, making tampering detectable by traversing the chain and verifying each link.

The Fact Layer enforces audit trail integrity through two complementary mechanisms. The first is the \textbf{state transition validator}, which intercepts every proposed agent action $a$ before it is committed to the ledger and validates it against the boundary specification $B$ sourced from the Purpose Layer. An action that violates any constraint in $B$ is rejected by the validator and replaced with a \texttt{bailout\underline{ }trigger} entry in the ledger—the protocol treats a rejected action as a bail-out event, not merely a denied action. This ensures that even blocked or prevented actions leave a forensic record. The second mechanism is the \textbf{checkpoint compulsion rule}: the Fact Layer requires that a checkpoint $c_t$ be produced and written to the ledger at every interval $\tau_{\text{ckpt}} \leq T_{\text{safety}}$, without exception. The absence of a checkpoint at the expected interval constitutes a ledger anomaly and is itself logged as a \texttt{checkpoint\underline{ }gap} event, triggering an automatic principal notification.

The Fact Layer's enforcement of the audit trail is what makes the Bailout Protocol's Accountability property (Property A) architecturally verifiable rather than institutionally trusted. The attribution chain $\phi_t$ within each ledger entry links every action to the principal $p$ who was authenticated at the time of the action's authorization. The chain is sealed by the principal's session credential, which is in turn linked to their authenticated identity in the Purpose Layer. The result is a complete, tamper-evident, time-stamped record of every decision and action in the session's lifecycle—from the initial Purpose Layer entry through to session termination—making post-hoc audit a matter of ledger traversal rather than reconstruction from scattered logs.

\subsection{Runtime Expression: Bailout Protocol as BX3's Accountability Guarantee in Operation}

The preceding subsections establish the structural components. Here we articulate the key relationship: the \textbf{Bailout Protocol is the runtime expression of BX3's upstream accountability guarantee}. This is not a metaphorical claim—it is an architectural one with practical consequences for protocol compliance and certification.

The BX3 Framework's accountability guarantee is established upstream, at the moment a human principal encodes their authorization in the Purpose Layer. When $p$ records $(p, B, J)$ at session initialization, they are making a formal commitment: \emph{within the bounds $B$, agent actions are authorized; outside those bounds, they are not}. This commitment is not self-executing—it does not prevent the agent from taking out-of-bounds actions, nor does it automatically halt the agent when bounds are approached. The guarantee becomes operative only when there exists a mechanism that (i) observes whether the agent's running state remains within $B$, (ii) detects the moment bounds are reached or approached, and (iii) executes a defined response—suspension, principal notification, or termination—when the bounds are violated.

The Bailout Protocol is precisely that mechanism. It is the operationalization of the upstream guarantee: the formal promise encoded in the Purpose Layer is transformed into a runtime monitoring, detection, and intervention protocol by the joint action of the Bounds Engine, the Fact Layer, and the state machine. Without the Bailout Protocol, the Purpose Layer's accountability guarantee would be a static declaration—a document rather than a system. The guarantee would hold only to the extent that the agent chose to honor it, which is precisely the condition the framework is designed to eliminate. BX3's accountability architecture is only as strong as its runtime enforcement; the Bailout Protocol is that enforcement.

This relationship has a direct implication for compliance verification. To verify that a BX3-based agentic system maintains its accountability guarantee, one need not examine the agent's decision logic directly—one examines whether the Bailout Protocol is correctly instantiated and operating. The audit trail in the Fact Layer provides the evidence. If the ledger is complete, the chain is intact, and the session history shows no periods in which the Bounds Engine was disabled or the Fact Layer's logging was suspended, then the upstream guarantee has been upheld. This inverts the traditional verification burden: rather than proving that the agent is trustworthy and then trusting its self-reports, BX3 treats the protocol itself as the primary evidence of compliance.

\subsection{State Machine as Fact Layer Extension}

The Bailout Protocol's session dynamics are formally captured by a \textbf{state machine} that is implemented as an extension of the Fact Layer. This state machine governs the session lifecycle, defines the legal state transitions, and records every transition event as a forensic ledger entry. The state machine is not a separate abstraction—it is the Fact Layer's operational controller, the mechanism by which the Fact Layer's logging regime is enforced and by which the Bounds Engine's trigger events are translated into system-level state changes.

The Bailout Protocol state machine comprises the following canonical states:

\begin{itemize}
    \item \texttt{boot}: The session is initializing. Purpose Layer entry $(p, B, J)$ is being loaded; the Bounds Engine is being parameterized with $B$; the forensic ledger is being opened with $J$ as the genesis entry.
    \item \texttt{active}: The agent is operating normally. Checkpoints are being produced at interval $\tau_{\text{ckpt}}$; the Bounds Engine is evaluating predicates; the Fact Layer is recording ledger entries.
    \item \texttt{bailout\underline{ }pending}: A boundary violation or timeout has been detected. The agent has been notified and is in a suspension state. The principal $p$ has been notified. The system awaits principal response.
    \item \texttt{escalation}: The principal has not responded within $T_{\text{liveness}}$. Registered fallback principals are being notified. The agent remains suspended.
    \item \texttt{recovery}: A principal has authenticated and is reviewing the session state. The session may be resumed on principal authorization or terminated based on principal directive.
    \item \texttt{terminated}: The session has ended. The forensic ledger is sealed (final hash computed, no further entries permitted). The session state is archived for the audit window $T_{\text{audit}}$.
\end{itemize}

The state machine defines a deterministic transition function $\delta: S \times E \rightarrow S$, where $S$ is the state set and $E$ is the event set. Events $E$ include: \texttt{init} (Purpose Layer loaded), \texttt{tick} (checkpoint interval elapsed), \texttt{boundary\underline{ }violation}, \texttt{timeout}, \texttt{principal\underline{ }auth} (human authenticated), \texttt{fallback\underline{ }notify}, \texttt{resume}, and \texttt{terminate}. Every transition $\delta(s, e)$ is logged to the forensic ledger as a \texttt{state\underline{ }transition} entry containing the source state, the event, the destination state, and the timestamp. This makes the session's entire lifecycle mechanically legible: any reviewer can reconstruct the session's evolution by traversing the ledger's state-transition entries.

The state machine's extension of the Fact Layer serves two purposes. First, it makes the Fact Layer's logging regime stateful: the ledger entries are not merely a chronological record but a mechanically interpretable trace of a defined automaton. Second, it grounds the Bailout Protocol's formal correctness theorem (Section~\ref{sec:guarantees}) in a concrete implementation artifact. The triadic invariant—Safety $\land$ Liveness $\land$ Accountability—is verified by model-checking the state machine: the invariant holds if and only if all reachable states from the initial state \texttt{boot} satisfy the invariant's conditions, and all transition paths that would violate the invariant are intercepted by a mandatory transition to \texttt{bailout\underline{ }pending} before the violation can occur.

The state machine is implemented in the Fact Layer's execution context, independent of the agent's process. State transitions are triggered by events sourced from the Bounds Engine (for \texttt{boundary\underline{ }violation} and \texttt{timeout}), from the authentication subsystem (for \texttt{principal\underline{ }auth}), and from the principal themselves (for \texttt{resume} and \texttt{terminate}). This decoupling ensures that the state machine's logic is not susceptible to manipulation by the agent it governs—a compromised agent cannot forge a state transition event, suppress a bailout trigger, or alter a ledger entry. The Fact Layer's append-only regime and cryptographic sealing apply to state machine events with the same rigor as to agent actions, ensuring that even the session's control flow history is tamper-evident.

\subsection{Forensic Ledger: Structural Properties}

The forensic ledger is the unifying data structure of the Bailout Protocol's BX3 implementation. It is written by the Fact Layer, seeded by the Purpose Layer, and read by the Bounds Engine for trajectory evaluation. Its structural properties ensure that it is both complete and honest—two properties that are separately necessary and jointly sufficient for it to serve as the definitive record of session accountability.

\begin{itemize}
    \item \textbf{Append-only.} No entry may be modified or deleted after writing. Edits would break the hash chain and be detectable by routine verification. The ledger's append-only property is enforced by the Fact Layer's write validator, which rejects any write that would alter a historical entry.
    \item \textbf{Complete.} Every checkpoint, every action, every state transition, every Bounds Engine predicate evaluation, every principal notification, and every human authentication event is recorded. The ledger contains no null intervals without a \texttt{checkpoint\underline{ }gap} entry recording the reason. Completeness is verified by the Fact Layer's checkpoint compulsion rule.
    \item \textbf{Tamper-evident.} The chained hash $\tau_t$ ensures that any post-hoc modification of a historical entry is detectable by traversing the chain and comparing the recomputed hash with the recorded hash. A mismatch indicates a breach of integrity.
    \item \textbf{Attributed.} Every entry contains a principal attribution reference $\phi_t$ linking it to the authenticated principal who was active at the time of the entry's creation. The attribution chain is non-repudiable: a principal cannot claim ignorance of an action whose attribution chain points to their authenticated session.
    \item \textbf{Bounded retention.} The ledger is retained for the configured audit window $T_{\text{audit}}$, after which it is cryptographically sealed and archived. The archive is write-once, read-many, with no mechanism for modification or deletion during the retention period. This satisfies regulatory and institutional requirements for audit record preservation.
\end{itemize}

These properties collectively establish the forensic ledger as the definitive record against which the Bailout Protocol's guarantees are verified. An auditor examining the ledger can determine, for any point in the session's history: who was the accountable principal, what bounds were in effect, what the agent's state was, what actions were taken, and whether any bailout event was triggered. The ledger is self-certifying: its integrity can be verified mechanically, without relying on the agent or any external testimony.
